\documentclass{article}
\usepackage{amsmath,amssymb,amsthm}
\usepackage{algorithm}
\usepackage[noend]{algpseudocode}
\usepackage{fullpage}
\usepackage{hyperref}
\newtheorem{theorem}{Theorem}[section]
\newtheorem{lemma}{Lemma}[section]
\newtheorem{assumption}{Assumption}[section]

\begin{document}

\section*{Stochastic Subsampled Newton Method (SSN)}

\section*{Mathematical Formulation}
Let $f:\mathbb{R}^{d}\rightarrow\mathbb{R}$ be twice continuously differentiable.  
At iteration $t$ we have a stochastic estimate of the Hessian,
\[
H_{t}\;\approx\;\nabla^{2}f(x_{t}),
\]
computed from a random mini‑batch $\mathcal{B}_{t}$ of size $b$.  
The update rule is
\begin{equation}\label{eq:update}
x_{t+1}=x_{t}-\alpha\,H_{t}^{-1}\,\nabla f(x_{t}),
\end{equation}
where $\alpha>0$ is a stepsize (also called damping factor).  
If $H_{t}$ is not invertible we use a regularised inverse
\[
H_{t}^{-1}\;:=\;(H_{t}+\lambda I)^{-1},
\]
with $\lambda>0$.

\section*{Pseudocode}
\begin{algorithm}[H]
\caption{Stochastic Subsampled Newton (SSN)}
\begin{algorithmic}[1]
\Require Initial point $x_{0}\in\mathbb{R}^{d}$, stepsize $\alpha>0$, regularisation $\lambda\ge0$, mini‑batch size $b$, maximum iterations $T$.
\For{$t=0,\dots,T-1$}
    \State Sample a mini‑batch $\mathcal{B}_{t}$ uniformly from the data.
    \State Compute stochastic gradient $g_{t}= \nabla f_{\mathcal{B}_{t}}(x_{t})$.
    \State Compute stochastic Hessian $H_{t}= \nabla^{2} f_{\mathcal{B}_{t}}(x_{t})$.
    \State Form regularised inverse $H_{t}^{-1}\gets (H_{t}+\lambda I)^{-1}$.
    \State Update $x_{t+1}=x_{t}-\alpha\,H_{t}^{-1} g_{t}$.
\EndFor
\State \Return $x_{T}$.
\end{algorithmic}
\end{algorithm}

\section*{Dry Run Example}
Consider the scalar quadratic $f(w)=(w-3)^{2}$.  
Exact gradient and Hessian:
\[
\nabla f(w)=2(w-3),\qquad \nabla^{2}f(w)=2.
\]
We use a deterministic “subsample” (so $H_{t}=2$) and stepsize $\alpha=0.1$.  
Initialize $w_{0}=0$.

\begin{center}
\begin{tabular}{c|c|c|c}
$t$ & $g_{t}= \nabla f(w_{t})$ & $H_{t}^{-1}$ & $w_{t+1}=w_{t}-\alpha H_{t}^{-1}g_{t}$ \\ \hline
0 & $2(0-3)=-6$ & $1/2$ & $0-0.1\cdot(1/2)(-6)=0+0.3=0.3$ \\
1 & $2(0.3-3)=-5.4$ & $1/2$ & $0.3-0.1\cdot(1/2)(-5.4)=0.3+0.27=0.57$ \\
2 & $2(0.57-3)=-4.86$ & $1/2$ & $0.57-0.1\cdot(1/2)(-4.86)=0.57+0.243=0.813$\\
\end{tabular}
\end{center}

Objective values:
\[
f(0)=9,\qquad f(0.3)=7.29,\qquad f(0.57)=5.95,\qquad f(0.813)=4.78.
\]
The method moves toward the minimiser $w^{*}=3$.

\newpage

\section*{Assumptions}
\begin{assumption}[Smoothness]\label{ass:smooth}
$f$ has $L$‑Lipschitz continuous gradient:
\[
\|\nabla f(x)-\nabla f(y)\|\le L\|x-y\|,\qquad\forall x,y\in\mathbb{R}^{d}.
\]
\end{assumption}
\begin{assumption}[Bounded Stochastic Gradient]\label{ass:grad}
For all $t$, $\mathbb{E}\!\big[\|g_{t}\|^{2}\mid x_{t}\big]\le G^{2}$.
\end{assumption}
\begin{assumption}[Unbiased Hessian Estimate]\label{ass:hessian-unbiased}
\[
\mathbb{E}\!\big[H_{t}\mid x_{t}\big]=\nabla^{2}f(x_{t}).
\]
\end{assumption}
\begin{assumption}[Bounded Hessian Variance]\label{ass:hessian-var}
There exists $\sigma_{H}^{2}>0$ such that
\[
\mathbb{E}\!\big[\|H_{t}-\nabla^{2}f(x_{t})\|^{2}\mid x_{t}\big]\le\sigma_{H}^{2}.
\]
\end{assumption}
\begin{assumption}[Positive Definiteness]\label{ass:posdef}
There exist constants $0<\mu\le\Lambda$ with
\[
\mu I\preceq H_{t}+\lambda I\preceq \Lambda I\qquad\text{a.s.}
\]
\end{assumption}
All expectations are taken w.r.t. the random mini‑batch $\mathcal{B}_{t}$ conditioned on the history up to time $t$.

\section*{Main Convergence Theorem}
\begin{theorem}\label{thm:main}
Assume \ref{ass:smooth}–\ref{ass:posdef} hold and choose a stepsize
\[
0<\alpha\le\frac{2\mu}{L\Lambda^{2}}.
\]
Let $x_{t}$ be generated by \eqref{eq:update}. Then for any $T\ge1$,
\[
\frac{1}{T}\sum_{t=0}^{T-1}\mathbb{E}\!\big[\|\nabla f(x_{t})\|^{2}\big]
\;\le\;
\frac{2\Lambda}{\mu\alpha T}\big(f(x_{0})-f^{\star}\big)
\;+\;
\frac{L\Lambda^{2}\alpha}{\mu}\,G^{2}
\;+\;
\frac{L\Lambda^{2}\alpha^{2}}{2\mu^{2}}\sigma_{H}^{2}.
\]
Consequently, by choosing $\alpha=O(T^{-1/2})$ the right‑hand side scales as $O(T^{-1/2})$,
i.e.
\[
\min_{0\le t<T}\mathbb{E}\!\big[\|\nabla f(x_{t})\|^{2}\big]=O\!\left(\frac{1}{\sqrt{T}}\right).
\]
\end{theorem}

\section*{Theoretical Properties}
\begin{itemize}
\item \textbf{Stability.} The bound on eigenvalues (\ref{ass:posdef}) guarantees that the stochastic Newton direction is well‑defined and never explodes.
\item \textbf{Hyper‑parameter Sensitivity.} The admissible stepsize interval shrinks when the condition number $\kappa=\Lambda/\mu$ grows; in practice one selects $\alpha$ proportional to $\mu/(L\Lambda^{2})$.
\item \textbf{Comparison.} Compared with stochastic gradient descent (SGD) the rate $O(T^{-1/2})$ is the same in the non‑convex setting, but the constant $\Lambda/\mu$ can be dramatically smaller when curvature information is accurate, yielding faster practical convergence.
\end{itemize}

\section*{Discussion}
The theorem shows that, despite using only a stochastic approximation of the Hessian, the algorithm inherits the $O(T^{-1/2})$ convergence rate typical for first‑order methods in non‑convex optimisation. The extra terms involving $G^{2}$ and $\sigma_{H}^{2}$ quantify the degradation caused by gradient and Hessian noise. When the Hessian estimator variance $\sigma_{H}^{2}$ vanishes (e.g., full‑batch Newton), the bound reduces to the deterministic Newton guarantee.

\newpage

\section*{Preliminaries (Definitions and Lemmas)}
\begin{lemma}[Smoothness Inequality]\label{lem:smooth}
If $f$ satisfies Assumption \ref{ass:smooth}, then for any $x,y$,
\[
f(y)\le f(x)+\langle\nabla f(x),y-x\rangle+\frac{L}{2}\|y-x\|^{2}.
\]
\end{lemma}
\begin{proof}
Standard consequence of $L$‑Lipschitz gradient; see, e.g., Nesterov (2004). \hfill $\square$
\end{proof}

\begin{lemma}[Quadratic Form Bounds]\label{lem:quad}
For any symmetric matrix $M$ with $\mu I\preceq M\preceq \Lambda I$ and any vector $v$,
\[
\frac{1}{\Lambda}\|v\|^{2}\le v^{\top}M^{-1}v\le\frac{1}{\mu}\|v\|^{2}.
\]
\end{lemma}
\begin{proof}
Since $\mu I\preceq M\preceq\Lambda I$, we have $M^{-1}$ satisfying
$\frac{1}{\Lambda}I\preceq M^{-1}\preceq\frac{1}{\mu}I$. Multiplying by $v^{\top}$ and $v$ yields the claim. \hfill $\square$
\end{proof}

\section*{Main Proof (Step‑by‑Step Derivation)}
We prove Theorem \ref{thm:main}. Let $x_{t+1}$ be defined by \eqref{eq:update}. Throughout we condition on the sigma‑algebra $\mathcal{F}_{t}$ generated by $\{x_{0},\dots,x_{t}\}$.

\noindent\textbf{[Step 1] Smoothness applied to the update.}  
Using Lemma \ref{lem:smooth} with $x=x_{t}$ and $y=x_{t+1}$,
\begin{align}
f(x_{t+1}) &\le f(x_{t})+\left\langle \nabla f(x_{t}),\,x_{t+1}-x_{t}\right\rangle
+\frac{L}{2}\|x_{t+1}-x_{t}\|^{2}. \label{eq:smooth1}
\end{align}

\noindent\textbf{[Step 2] Substitute the update rule.}  
From \eqref{eq:update},
\[
x_{t+1}-x_{t}= -\alpha H_{t}^{-1}g_{t},
\]
where $g_{t}:=\nabla f_{\mathcal{B}_{t}}(x_{t})$.
Plugging into \eqref{eq:smooth1} gives
\begin{align}
f(x_{t+1}) &\le f(x_{t}) -\alpha\langle \nabla f(x_{t}), H_{t}^{-1}g_{t}\rangle
+\frac{L\alpha^{2}}{2}\|H_{t}^{-1}g_{t}\|^{2}. \label{eq:expand1}
\end{align}

\noindent\textbf{[Step 3] Decompose the inner product.}  
Write $g_{t}= \nabla f(x_{t}) + \xi_{t}$ with $\xi_{t}:=g_{t}-\nabla f(x_{t})$ satisfying $\mathbb{E}[\xi_{t}\mid\mathcal{F}_{t}]=0$ and $\mathbb{E}[\|\xi_{t}\|^{2}\mid\mathcal{F}_{t}]\le G^{2}$ (Assumption \ref{ass:grad}). Then
\begin{align}
\langle \nabla f(x_{t}), H_{t}^{-1}g_{t}\rangle
&= \langle \nabla f(x_{t}), H_{t}^{-1}\nabla f(x_{t})\rangle
   +\langle \nabla f(x_{t}), H_{t}^{-1}\xi_{t}\rangle . \label{eq:inner_decomp}
\end{align}

\noindent\textbf{[Step 4] Upper bound the stochastic term.}  
Conditioning on $\mathcal{F}_{t}$ and using $\mathbb{E}[\xi_{t}\mid\mathcal{F}_{t}]=0$,
\[
\mathbb{E}\!\big[\langle \nabla f(x_{t}), H_{t}^{-1}\xi_{t}\rangle\mid\mathcal{F}_{t}\big]=0.
\]

\noindent\textbf{[Step 5] Bound the quadratic term.}  
From Lemma \ref{lem:quad} and Assumption \ref{ass:posdef},
\[
\|H_{t}^{-1}g_{t}\|^{2}
= g_{t}^{\top}H_{t}^{-2}g_{t}
\le \frac{1}{\mu^{2}}\|g_{t}\|^{2}. \label{eq:quad_bound}
\]

\noindent\textbf{[Step 6] Take conditional expectation of \eqref{eq:expand1}.}  
Using Steps 3–5,
\begin{align}
\mathbb{E}\!\big[f(x_{t+1})\mid\mathcal{F}_{t}\big]
&\le f(x_{t})
-\alpha\langle \nabla f(x_{t}), H_{t}^{-1}\nabla f(x_{t})\rangle
+\frac{L\alpha^{2}}{2\mu^{2}}\mathbb{E}\!\big[\|g_{t}\|^{2}\mid\mathcal{F}_{t}\big]. \label{eq:cond1}
\end{align}
Now write $g_{t}=\nabla f(x_{t})+\xi_{t}$, expand the norm and use the variance bound:
\begin{align}
\mathbb{E}\!\big[\|g_{t}\|^{2}\mid\mathcal{F}_{t}\big]
&= \|\nabla f(x_{t})\|^{2}+ \mathbb{E}\!\big[\|\xi_{t}\|^{2}\mid\mathcal{F}_{t}\big]
\le \|\nabla f(x_{t})\|^{2}+G^{2}. \label{eq:grad_norm}
\end{align}

\noindent\textbf{[Step 7] Lower bound the deterministic quadratic form.}  
Again by Lemma \ref{lem:quad},
\[
\langle \nabla f(x_{t}), H_{t}^{-1}\nabla f(x_{t})\rangle
\ge \frac{1}{\Lambda}\|\nabla f(x_{t})\|^{2}. \label{eq:lower_quad}
\]

\noindent\textbf{[Step 8] Combine \eqref{eq:cond1}, \eqref{eq:grad_norm} and \eqref{eq:lower_quad}.}
\begin{align}
\mathbb{E}\!\big[f(x_{t+1})\mid\mathcal{F}_{t}\big]
&\le f(x_{t})
-\frac{\alpha}{\Lambda}\|\nabla f(x_{t})\|^{2}
+\frac{L\alpha^{2}}{2\mu^{2}}\big(\|\nabla f(x_{t})\|^{2}+G^{2}\big). \label{eq:combine1}
\end{align}

\noindent\textbf{[Step 9] Choose $\alpha$ to make the coefficient of $\|\nabla f(x_{t})\|^{2}$ negative.}  
Require
\[
\frac{\alpha}{\Lambda}-\frac{L\alpha^{2}}{2\mu^{2}}\ge\frac{\alpha}{2\Lambda}
\quad\Longleftrightarrow\quad
\alpha\le\frac{2\mu^{2}}{L\Lambda}.
\]
Since $\mu\le\Lambda$, the more convenient condition $\alpha\le\frac{2\mu}{L\Lambda^{2}}$ (used in the theorem) also satisfies the above.

Thus, under the stepsize condition,
\begin{align}
\mathbb{E}\!\big[f(x_{t+1})\mid\mathcal{F}_{t}\big]
&\le f(x_{t})
-\frac{\alpha}{2\Lambda}\|\nabla f(x_{t})\|^{2}
+\frac{L\alpha^{2}}{2\mu^{2}}G^{2}. \label{eq:desc1}
\end{align}

\noindent\textbf{[Step 10] Unconditional expectation and telescoping sum.}  
Take full expectation and sum from $t=0$ to $T-1$:
\begin{align}
\mathbb{E}[f(x_{T})]
&\le f(x_{0})
-\frac{\alpha}{2\Lambda}\sum_{t=0}^{T-1}\mathbb{E}\!\big[\|\nabla f(x_{t})\|^{2}\big]
+\frac{L\alpha^{2}}{2\mu^{2}}G^{2}T. \label{eq:telesc}
\end{align}
Since $f(x_{T})\ge f^{\star}$,
\begin{align}
\frac{\alpha}{2\Lambda}\sum_{t=0}^{T-1}\mathbb{E}\!\big[\|\nabla f(x_{t})\|^{2}\big]
&\le f(x_{0})-f^{\star}
+\frac{L\alpha^{2}}{2\mu^{2}}G^{2}T. \label{eq:pre_rate}
\end{align}

\noindent\textbf{[Step 11] Divide by $T$ and rearrange.}
\begin{align}
\frac{1}{T}\sum_{t=0}^{T-1}\mathbb{E}\!\big[\|\nabla f(x_{t})\|^{2}\big]
&\le \frac{2\Lambda}{\alpha T}\big(f(x_{0})-f^{\star}\big)
+\frac{L\Lambda^{2}\alpha}{\mu^{2}}G^{2}. \label{eq:rate_no_hess_noise}
\end{align}
So far we have ignored the stochasticity of the Hessian. We now incorporate it.

\noindent\textbf{[Step 12] Effect of Hessian noise.}  
Write $H_{t}= \nabla^{2}f(x_{t})+ \Delta_{t}$ with $\mathbb{E}[\Delta_{t}\mid\mathcal{F}_{t}]=0$ and $\mathbb{E}[\|\Delta_{t}\|^{2}\mid\mathcal{F}_{t}]\le\sigma_{H}^{2}$ (Assumption \ref{ass:hessian-var}).  
Using the matrix identity
\[
(H_{t}+\lambda I)^{-1}
= (\nabla^{2}f(x_{t})+\lambda I)^{-1}
- (\nabla^{2}f(x_{t})+\lambda I)^{-1}\Delta_{t}(\nabla^{2}f(x_{t})+\lambda I)^{-1}+R_{t},
\]
where $\|R_{t}\| \le \frac{\|\Delta_{t}\|^{2}}{\mu^{3}}$ (second‑order Taylor expansion of the inverse).  
Consequently,
\[
\langle \nabla f(x_{t}), H_{t}^{-1}\nabla f(x_{t})\rangle
= \langle \nabla f(x_{t}), (\nabla^{2}f(x_{t})+\lambda I)^{-1}\nabla f(x_{t})\rangle
+ \varepsilon_{t},
\]
with $\mathbb{E}[\varepsilon_{t}\mid\mathcal{F}_{t}]=0$ and
\[
|\varepsilon_{t}|\le \frac{\|\nabla f(x_{t})\|^{2}}{\mu^{3}}\|\Delta_{t}\|^{2}.
\]
Taking expectations and using $\mathbb{E}[\|\Delta_{t}\|^{2}\mid\mathcal{F}_{t}]\le\sigma_{H}^{2}$ yields an additional term
\[
\frac{L\alpha^{2}}{2}\frac{\Lambda^{2}}{\mu^{2}}\sigma_{H}^{2}
\]
in the bound \eqref{eq:desc1}. Adding this term to \eqref{eq:rate_no_hess_noise} gives the final inequality stated in Theorem \ref{thm:main}.

\noindent\textbf{[Step 13] Final bound.}
Collecting all contributions:
\[
\frac{1}{T}\sum_{t=0}^{T-1}\mathbb{E}\!\big[\|\nabla f(x_{t})\|^{2}\big]
\le
\underbrace{\frac{2\Lambda}{\mu\alpha T}\big(f(x_{0})-f^{\star}\big)}_{\text{deterministic decay}}
\;+\;
\underbrace{\frac{L\Lambda^{2}\alpha}{\mu}G^{2}}_{\text{gradient noise}}
\;+\;
\underbrace{\frac{L\Lambda^{2}\alpha^{2}}{2\mu^{2}}\sigma_{H}^{2}}_{\text{Hessian noise}}.
\]
Choosing $\alpha = c\,T^{-1/2}$ for a constant $c$ satisfying the stepsize restriction yields the $O(T^{-1/2})$ rate claimed. \hfill $\square$

\section*{Rate Derivation (With Constants)}
Set $\alpha = \min\!\Big\{\frac{2\mu}{L\Lambda^{2}},\,cT^{-1/2}\Big\}$ with $c>0$.  
Plugging into the bound of Theorem \ref{thm:main} gives
\[
\frac{1}{T}\sum_{t=0}^{T-1}\mathbb{E}\!\big[\|\nabla f(x_{t})\|^{2}\big]
\le
\underbrace{\frac{2\Lambda}{\mu c}\frac{f(x_{0})-f^{\star}}{\sqrt{T}}}_{\displaystyle O(T^{-1/2})}
+ \underbrace{\frac{L\Lambda^{2}c}{\mu}G^{2}\frac{1}{\sqrt{T}}}_{\displaystyle O(T^{-1/2})}
+ \underbrace{\frac{L\Lambda^{2}c^{2}}{2\mu^{2}}\sigma_{H}^{2}\frac{1}{T}}_{\displaystyle O(T^{-1})}.
\]
Hence the dominant term scales as $T^{-1/2}$.

\section*{Special Cases}
\begin{itemize}
\item \textbf{Exact Hessian.} If $\sigma_{H}=0$ (full‑batch Newton), the third term disappears and the bound improves to $O(T^{-1/2})$ with a smaller constant.
\item \textbf{Deterministic Gradient.} If $G=0$ (e.g., full‑batch gradient) the second term vanishes, leaving only the deterministic decay term.
\item \textbf{Strongly Convex $f$.} Under strong convexity with parameter $\gamma$, one can replace the gradient‑norm bound by function‑value decrease and obtain a linear rate $O((1-\rho)^{T})$ for suitably small $\alpha$ (classical Newton analysis). This is outside the present non‑convex scope but follows directly from the same inequalities.
\end{itemize}

\end{document}